\section{Trabajo relacionado}
\label{sec:relacionado}

Organizamos el trabajo relacionado por lo que cada línea establece y por dónde se detiene, porque
el hueco que este artículo aborda está en el espacio entre dos literaturas que rara vez se
encuentran.

\paragraph{Clasificación selectiva.}
La opción de abstenerse es antigua y está bien caracterizada. La regla de Chow fija el compromiso
óptimo entre error y rechazo para una posterior conocida \cite{chow1970reject}, y el tratamiento
moderno la extiende a redes profundas con rechazo aprendido o calibrado
\cite{geifman2017selective,geifman2019selectivenet,liu2019deepgamblers}, con revisiones recientes
que cartografían el espacio de diseño \cite{hendrickx2024rejectsurvey,hasan2023uncertaintysurvey}.
Dos refinamientos importan aquí: la observación de que una sola curva de riesgo y cobertura oculta
qué grupos pagan la abstención \cite{jones2020selectivedisparities}, y la extensión de las curvas
de rechazo a precisión y exhaustividad en lugar de solo al error
\cite{fischer2023rejectcurves,pugnana2022aucselective}. Lo que esta literatura no trata es la
abstención \emph{por clase}: su unidad de rechazo es la muestra, así que un mecanismo que declina
nombrar una categoría entera cae fuera de su formalismo, y la interacción entre ese mecanismo y
una métrica promediada por clase no se analiza.

\paragraph{Estructura de clases en mapeo de cultivos.}
El trabajo sobre mapeo de cultivos sí manipula el catálogo de clases, pero como recurso de
modelado y no como punto de operación. Turkoglu y colaboradores organizan los tipos de cultivo en
una jerarquía y dejan que el modelo retroceda a un nivel más grueso cuando la decisión fina no es
fiable \cite{turkoglu2021hierarchies}, que es el mecanismo existente más cercano al recorte de
leyenda; la evaluación, sin embargo, reporta exactitud global sobre una sola partición, así que ni
el efecto sobre un promedio por clase ni su variabilidad en el espacio quedan medidos. Los
esfuerzos de armonización paneuropea afrontan el mismo problema al nivel de la nomenclatura
\cite{schneider2021eurocrops,claverie2026eurocropsv2}, y el trabajo de mapeo de grano fino señala
el costo de las clases raras sin aislarlo
\cite{lei2025finecrop,li2021costeffective,donmez2025cropcover}. En toda esta literatura el catálogo
podado se reporta como una propiedad del dato.

\paragraph{Bancos de datos y modelos temporales.}
Evaluamos sobre dos bancos públicos por parcela. PASTIS aporta series Sentinel-2 con anotación
panóptica sobre Francia \cite{garnot2021pastis}; BreizhCrops aporta una colección por parcela más
grande sobre Bretaña, con partición nativa por región administrativa y un reparto de clases extremo
\cite{russwurm2020breizhcrops}. Las arquitecturas temporales que usamos como predictores son
estándar: convoluciones temporales \cite{pelletier2019tempcnn}, autoatención sobre series crudas
\cite{russwurm2020selfattention} y transformadores de visión adaptados a series de imágenes
satelitales \cite{tarasiou2023tsvit}. Ninguna es contribución de este artículo; aportan las
posteriores sobre las que opera el protocolo.

\paragraph{Protocolo de evaluación.}
Que el protocolo sea espacial y no aleatorio no es opcional. La validación cruzada que ignora la
autocorrelación espacial infla las estimaciones de capacidad predictiva, a veces de forma severa
\cite{roberts2017crossvalidation,ploton2020spatialvalidation}. Por eso elegimos cada punto de
operación sobre bloques espaciales disjuntos del bloque donde se evalúa. La literatura de ensambles
aporta el ingrediente restante, el meta-modelo de apilamiento cuyo comportamiento implícito
examinamos \cite{wolpert1992stacking,zhang2022ensemblereview,abdali2023parallelcascaded}, y es
precisamente la exigencia de Wolpert de particiones libres de fuga para entrenar el modelo de
segundo nivel la que nuestra comparación entre dentro y fuera de muestra pone a prueba.
